% ============================================================================================
\section{Data Handling and Preprocessing}
\label{sec:data}
Every decision below is also logged, with counts, in \fpath{notes/preprocessing_log.md} (reproduced in
Appendix~\ref{app:log}).

\subsection{gridMET download and inspection}
\textbf{Download.} \fpath{download_gridmet.py} fetches \texttt{vs, th, tmmx, pr, vpd, rmax, srad} for 2024
(28--205\,MB each, about 820\,MB in total) from the repository listing, which we verified rather than assumed. It shows
a progress bar and resumes partial downloads with HTTP \texttt{Range} requests; one interrupted download was resumed
successfully. Because the files are small, server-side OPeNDAP subsetting was not needed.

\textbf{Raw schema.} \fpath{inspect_raw.py} prints each file without CF decoding, so we see what is on disk.
Every variable is stored as \texttt{uint16} with \texttt{\_FillValue}=32767, under a long CF name that differs from the
file abbreviation. The names are \texttt{air\_temperature}, \texttt{precipitation\_amount}, \texttt{wind\_speed},
\texttt{wind\_from\_direction}, \texttt{mean\_vapor\_pressure\_deficit}, \texttt{relative\_humidity} and
\texttt{surface\_downwelling\_shortwave\_flux\_in\_air}. The packing is variable-specific:
temperature uses \texttt{scale\_factor}=0.1 and \texttt{add\_offset}=220\,K; VPD uses 0.01; direction uses 1.0; the
others use 0.1. Dimensions are \texttt{(day, lat, lon)} = (366, 585, 1386), \texttt{day} is ``days since 1900-01-01'',
and \textbf{latitude is stored in descending order} (49.40\,$\rightarrow$\,25.07\,$^\circ$N).

\subsection{Subsetting, masking and unit conversion}
\textbf{Lazy loading.} Files are opened with dask chunks equal to the on-disk chunking $(61, 98, 231)$, and one variable
is processed at a time, so the full year (594\,MB per variable when decoded) is never in memory at once.

\textbf{Temporal subset.} 1~Jul--30~Sep~2024 gives 92 days. \texttt{day} is renamed \texttt{time}.

\textbf{Orientation.} We \emph{sort latitude ascending} once during preprocessing. Every downstream routine then uses
the same convention: image origin at the bottom, \texttt{streamplot} requiring increasing $y$, and Marching-Squares row
0 being the southern edge. This removes a whole class of flipped-map errors.

\textbf{Masking.} CF decoding applies the scale and offset and turns the fill value into NaN. The land mask is
\emph{identical} for all seven variables and constant across all 92 days: 483,879 of 810,810 cells, checked cell by cell.
NaN cells (ocean, Canada, Mexico) are never interpolated. Plots show them in a neutral blue-grey labelled ``No data''.

\textbf{Units.} Temperature is converted with
\begin{equation}
T[\si{\celsius}] = 0.1\,n + 220 - 273.15,
\end{equation}
where $n$ is the stored integer. Precipitation stays in mm\,day$^{-1}$, wind in m\,s$^{-1}$, VPD in kPa, $r_{\max}$ in \%
and $s_{rad}$ in W\,m$^{-2}$. Outputs are float32 NetCDF with zlib compression
(\fpath{data/processed/gridmet_<var>_2024JAS.nc}).

\textbf{Plausibility of extremes.} No values were clipped, because all extremes have a physical explanation. The maximum
$\Tmax$ is 53.85\,\degC{} (6~Jul, Death Valley). The \emph{minimum} daily maximum is $-6.35$\,\degC{} (23~Aug,
41.40\,$^\circ$N 122.18\,$^\circ$W), which is the summit of Mt~Shasta at about 4,300\,m, so the 4\,km grid resolves
high-elevation cold spots. The largest daily rain total is 401.4\,mm (4~Aug, 27.40\,$^\circ$N 82.39\,$^\circ$W, Sarasota--Manatee
coast) in Debby's rain band.

\subsection{Wind decomposition}
gridMET stores the \emph{meteorological} direction $\theta$: where the wind blows \emph{from}, in degrees clockwise from
north. Quiver plots and streamlines need the vector the air moves \emph{toward}, so
\begin{equation}
u = -v_s\sin\theta,\qquad v = -v_s\cos\theta \qquad(\theta\ \text{in radians}).
\label{eq:wind}
\end{equation}
The built-in sanity check gives the following. A westerly ($\theta=270^\circ$) at 10\,m/s gives $u=+10$, $v=0$ (air
moving east). A northerly ($\theta=0^\circ$) gives $v=-10$. $\theta=90^\circ$ gives $u=-10$, and $\theta=180^\circ$ gives
$v=+10$. Values $\theta=0$ (63,867 cell-days) and $\theta=360$ (42,580) are equivalent and decompose identically. Only 39
cell-days have $v_s<0.5$\,m/s, so undefined calm-wind directions are negligible. Over Florida on D7 this gives
$(u,v)\approx(-6,+11)$\,m/s at Tampa from $\theta=151^\circ$, $v_s=12.5$\,m/s. That is a south-south-easterly wind blowing
toward the north-north-west, consistent with Helene's counter-clockwise circulation east of its centre. The u/v fields are
stored in \fpath{gridmet_wind_2024JAS.nc}. Cartopy then rotates the geographic east/north components into the Lambert
conformal map frame.

\subsection{Statistics and derived fields}
For every variable and day we store $n$, min, max, mean, standard deviation and the 1st/5th/25th/50th/75th/95th/99th
percentiles (\fpath{stats_per_day.csv}), plus the same over the whole period (Table~\ref{tab:global}).
Table~\ref{tab:perdate} lists the values on the eight dates. The anomaly used for diverging maps is
\begin{equation}
a_{ij}(t) = \Tmax{}_{ij}(t) - \frac{1}{92}\sum_{t'}\Tmax{}_{ij}(t'),
\label{eq:anom}
\end{equation}
so it is relative to each cell's own Jul--Sep 2024 mean. This removes the elevation and latitude signal and leaves only
weather.

\input{generated/tab_global}
\input{generated/tab_perdate}

\textbf{Normalisation formulas used in S1--S3.} A colour map $C:[0,1]\rightarrow\mathrm{RGB}$ is applied to a normalised
value $\hat{x}$. We use linear normalisation $\hat{x}=(x-x_{\min})/(x_{\max}-x_{\min})$ with either \emph{global} bounds
(over all 92 days) or \emph{local} bounds (the current day). Logarithmic normalisation is
$\hat{x}=(\log x-\log x_0)/(\log x_1-\log x_0)$ with $x_0=0.1$\,mm, applied only to $x>0$. The two-slope diverging
normalisation maps $[x_{\min},0]\rightarrow[0,0.5]$ and $[0,x_{\max}]\rightarrow[0.5,1]$, so zero always sits at the
neutral centre colour. Discrete classes use a \texttt{BoundaryNorm} with breaks $b_0<\dots<b_K$ and
$\hat{x}=k/(K-1)$ for $b_k\le x<b_{k+1}$.

\subsection{A1 cleaning}
\label{sec:a1clean}
\fpath{a1_preprocess.py} re-implements the A1 cleaning so that A2 is self-contained.
\begin{itemize}
\item Merge the 50 files (49,266 rows) and confirm there are no duplicate ids.
\item Parse both timestamp formats; the range is 27~Sep~2011 to 11~Sep~2024 UTC.
\item Drop five columns: \texttt{num\_awards} (all zero), \texttt{crosspost\_subreddits} (99.35\,\% missing),
  \texttt{is\_bot} (8.53\,\% missing), \texttt{is\_megathread} (21 True) and \texttt{body} (82.79\,\% missing).
\item \textbf{Drop no rows.} The corpus is the top tail by construction, so extreme scores are the object of study and
  are handled with log scales.
\end{itemize}
New variables are year, month, UTC hour, decimal year, title length, $\log_{10}$ score and a five-level era
(2011--16, 2017--18, 2019--20, 2021--22, 2023--24, chosen so the sparse early years are pooled).

\textbf{Domain normalisation.} Domain strings are lower-cased, mobile and \texttt{www} prefixes are stripped, platform
aliases are merged (\texttt{youtu.be}$\rightarrow$\texttt{youtube.com}, \texttt{i.imgur.com}$\rightarrow$\texttt{imgur.com},
Reddit's \texttt{i./v./preview.redd.it}, \texttt{reddituploads.com} and \texttt{redditmedia.com}$\rightarrow$\texttt{redd.it}),
and each host is reduced to its registrable domain. For example \texttt{cdna.artstation.com} becomes
\texttt{artstation.com}, and \texttt{bbc.co.uk} keeps its second-level suffix. This took 3,420 raw strings down to 3,082
domains. Without it, one platform appeared as several unconnected nodes; it was found when the first network had
\texttt{cdna} and \texttt{cdnb.artstation.com} as separate nodes. Each post is then classified as a self post,
Reddit-hosted media (15,079), imgur (4,112), YouTube (3,040), other video host (890) or external website (13,010).

\textbf{Topic categories.} We grouped the 50 communities by purpose into eight topics: Q\&A and advice (8),
Text humour and stories (4), News, science and tech (8), Memes and GIFs (4), Photos and curiosities (7), Creative and
lifestyle (5), Entertainment and media (10), and Knowledge and curiosity (4). The topics are the common colour key in
I1--I3.

\subsection{Derived InfoVis structures}
\label{sec:derive}
\textbf{(1) Network for I1.} The graph is bipartite: subreddits $\leftrightarrow$ external domains. Edge $(s,d)$ exists if
subreddit $s$ has a top post linking to $d$, and its weight $w_{sd}$ is the number of such posts. Self posts and Reddit's
own media host are excluded because they say nothing about where a community links \emph{out} to. Domains with fewer
than three top posts are removed (3,035 $\rightarrow$ 628 domains) to suppress one-off links. The result has 671 nodes
(43 subreddits, 628 domains), 2,350 edges and total weight 18,245 posts, which is well beyond a toy graph. The seven
text-only communities (AskReddit, LifeProTips, NoStupidQuestions, Showerthoughts, explainlikeimfive, personalfinance,
relationship\_advice) have no qualifying link. Their absence is itself a finding. The more obvious crosspost graph was
rejected because \texttt{crosspost\_subreddits} is 99.35\,\% missing.

\textbf{(2) Hierarchies for I2b.} Algorithm~\ref{alg:tree} nests the flat table by an ordered list of categorical
levels. Changing the order or the levels \emph{remodels} the tree:
H1 topic$\rightarrow$subreddit$\rightarrow$post type (135 leaves);
H2 post type$\rightarrow$hosting class$\rightarrow$topic$\rightarrow$subreddit (250 leaves);
H3 era$\rightarrow$year$\rightarrow$topic$\rightarrow$subreddit (486 leaves).
Each leaf stores three additive measures (posts, total score, total comments), so the size encoding can also be
switched. That gives $3\times3$ distinct trees.

\begin{algorithm}[!t]
\caption{Tabular-to-tree remodel (\texttt{to\_tree})}
\label{alg:tree}
\begin{algorithmic}[1]
\REQUIRE table $R$, ordered level list $\mathcal{L}=(\ell_1,\dots,\ell_m)$, measures $\mathcal{M}$
\ENSURE nested node $\{name, children\}$
\STATE \textbf{function} \textsc{Rec}$(S, k, label)$
\IF{$k = m$}
  \RETURN $\{name{:}\,label\} \cup \{\mu{:}\sum_{r\in S} r.\mu \mid \mu\in\mathcal{M}\}$
\ENDIF
\STATE $C \leftarrow [\,]$
\FORALL{distinct value $g$ of $\ell_{k+1}$ in $S$ (sorted)}
  \STATE append \textsc{Rec}$(\{r\in S : r.\ell_{k+1}=g\}, k+1, g)$ to $C$
\ENDFOR
\RETURN $\{name{:}\,label,\ children{:}\,C\}$
\STATE \textbf{return} \textsc{Rec}$(R, 0, \text{root})$
\end{algorithmic}
\end{algorithm}

\textbf{(3) Multivariate table for I2a.} A stratified random sample of 40 posts per subreddit (seed 42) gives 2,000
records, which is enough to show structure while keeping interactive SVG brushing fast. It has seven quantitative axes
(decimal year, UTC hour, title length, upvote ratio, score, comments, crossposts) and one ordinal axis (post type). The
grouping variables are topic, post type and a k-means~\cite{lloyd1982} cluster. k-means runs on min--max-normalised axes,
with $\log_{10}$ applied to score and $\log_{10}(x+1)$ to comments and crossposts, the same transforms as the axis
scales. $k=4$ maximised the silhouette~\cite{rousseeuw1987} over $k=3..8$ at 0.248, with values of 0.241, 0.248, 0.229,
0.221, 0.226 and 0.194 for $k=3..8$. The low value shows that the clusters are soft, which is expected for engagement data,
and we report it rather than hide it. The cluster sizes are 552, 675, 553 and 220.

\textbf{(4) Time series for I3.} We build monthly counts of top posts and the monthly median $\log_{10}$(score), both per
topic and per subreddit, for Jan~2016--Aug~2024 (104 months, 98.3\,\% of posts). Before 2016 there are fewer than 50 posts
per month in total, and September~2024 is cut off at 11~Sep, so neither is used. Quarterly and yearly aggregates provide
the temporal-granularity comparison. Months with no top post are true zeros for counts, and are left blank for medians.
